% 第6章 线性空间与线性变换（同济《线性代数》第七版 第6章，选学但保留）
% §1 线性空间的定义与性质；§2 维数、基与坐标；§3 基变换与坐标变换；
% §4 线性变换；§5 线性变换的矩阵表示式。正文不写例题与习题。

\section{线性空间的定义与性质}

第4章的向量是"有序数组"，本章把向量与向量空间的概念公理化：元素是什么并不重要，重要的是元素之间定义了满足八条规律的加法与数乘。

\subsection{线性空间的定义}
\begin{definition}[线性空间]
设$V$是一个非空集合，$F$是一个数域（本书取实数域$\mathbb{R}$）。如果在$V$中定义了一个\textbf{加法}，即对任意$\boldsymbol{\alpha},\boldsymbol{\beta}\in V$，总有唯一的$\boldsymbol{\gamma}\in V$与之对应，记作$\boldsymbol{\gamma}=\boldsymbol{\alpha}+\boldsymbol{\beta}$；又在$V$中定义了一个\textbf{数与元素的乘法}（简称数乘），即对任一$\lambda\in F$与任一$\boldsymbol{\alpha}\in V$，总有唯一的$\boldsymbol{\delta}\in V$与之对应，记作$\boldsymbol{\delta}=\lambda\boldsymbol{\alpha}$；并且这两种运算满足以下八条运算规律（设$\boldsymbol{\alpha},\boldsymbol{\beta},\boldsymbol{\gamma}\in V$，$\lambda,\mu\in F$）：
\begin{enumerate}
    \item $\boldsymbol{\alpha}+\boldsymbol{\beta}=\boldsymbol{\beta}+\boldsymbol{\alpha}$；
    \item $(\boldsymbol{\alpha}+\boldsymbol{\beta})+\boldsymbol{\gamma}=\boldsymbol{\alpha}+(\boldsymbol{\beta}+\boldsymbol{\gamma})$；
    \item 在$V$中存在零元素$\boldsymbol{0}$，对任何$\boldsymbol{\alpha}\in V$都有$\boldsymbol{\alpha}+\boldsymbol{0}=\boldsymbol{\alpha}$；
    \item 对任何$\boldsymbol{\alpha}\in V$，存在$\boldsymbol{\alpha}$的负元素$\boldsymbol{\beta}\in V$，使$\boldsymbol{\alpha}+\boldsymbol{\beta}=\boldsymbol{0}$；
    \item $1\boldsymbol{\alpha}=\boldsymbol{\alpha}$；
    \item $\lambda(\mu\boldsymbol{\alpha})=(\lambda\mu)\boldsymbol{\alpha}$；
    \item $(\lambda+\mu)\boldsymbol{\alpha}=\lambda\boldsymbol{\alpha}+\mu\boldsymbol{\alpha}$；
    \item $\lambda(\boldsymbol{\alpha}+\boldsymbol{\beta})=\lambda\boldsymbol{\alpha}+\lambda\boldsymbol{\beta}$；
\end{enumerate}
那么就称$V$为（数域$F$上的）\textbf{线性空间}（或称向量空间），$V$中的元素统称为\textbf{向量}。
\end{definition}

八条规律的含义：$(i)(ii)$是加法的交换律、结合律；$(iii)(iv)$保证加法有逆运算；$(v)(vi)$与$(vii)(viii)$是数乘的结合律与分配律，$(v)$保证非零数乘有逆运算。满足这八条规律的加法与数乘统称为\textbf{线性运算}。

\corepoint{
线性空间$=$集合$+$线性运算。同一个集合若定义两种不同的线性运算，就构成不同的线性空间；若所定义的运算不满足八条规律，就不是线性空间。所以"运算"才是线性空间的本质，元素本身是什么并不重要。
}

常见与常见反例：
\begin{enumerate}
    \item 次数不超过$n$的多项式全体$P[x]_n=\{a_nx^{n}+\dots+a_1x+a_0\}$，按通常的多项式加法与数乘，构成线性空间（对运算封闭，八条规律自动满足）。
    \item 次数\textbf{恰等于}$n$的多项式全体，按通常运算\textbf{不}构成线性空间：两个$n$次多项式相加，最高次项可能抵消而次数降低，对加法不封闭。
    \item 形如$A\sin(x+B)$的函数全体，按函数加法与数乘构成线性空间。
    \item 正实数全体$\mathbb{R}^{+}$，若定义$a\oplus b=ab$，$\lambda\odot a=a^{\lambda}$，则构成线性空间（零元素是数$1$，$a$的负元素是$a^{-1}$）——这说明"向量"完全可以是数。
    \item 检验一个集合是否为线性空间时，若运算不是通常的实数加、乘，必须逐条验证八条规律，只验封闭性不够。
\end{enumerate}

\subsection{线性空间的性质}
设$V$是线性空间，则
\begin{enumerate}
    \item $V$中的零元素是唯一的；
    \item $V$中每个元素的负元素是唯一的，$\boldsymbol{\alpha}$的负元素记作$-\boldsymbol{\alpha}$；
    \item $0\boldsymbol{\alpha}=\boldsymbol{0}$，$(-1)\boldsymbol{\alpha}=-\boldsymbol{\alpha}$，$\lambda\boldsymbol{0}=\boldsymbol{0}$；
    \item 若$\lambda\boldsymbol{\alpha}=\boldsymbol{0}$，则$\lambda=0$或$\boldsymbol{\alpha}=\boldsymbol{0}$。
\end{enumerate}
（第4条在一般线性空间中不再"显然"，须由八条规律推出：若$\lambda\ne0$，两端乘$\lambda^{-1}$得$\boldsymbol{\alpha}=\lambda^{-1}\boldsymbol{0}=\boldsymbol{0}$。）

\subsection{子空间}
\begin{definition}[子空间]
设$V$是一个线性空间，$L$是$V$的一个非空子集，如果$L$对于$V$中所定义的加法与数乘两种运算也构成一个线性空间，则称$L$为$V$的\textbf{线性子空间}，简称\textbf{子空间}。
\end{definition}

\begin{theorem}[子空间的判定]
线性空间$V$的非空子集$L$构成$V$的子空间的充分必要条件是：$L$对于$V$中的线性运算封闭，即
\[
\boldsymbol{\alpha},\boldsymbol{\beta}\in L\ \Rightarrow\ \boldsymbol{\alpha}+\boldsymbol{\beta}\in L;
\qquad
\boldsymbol{\alpha}\in L,\ \lambda\in F\ \Rightarrow\ \lambda\boldsymbol{\alpha}\in L .
\]
\end{theorem}

因为$L\subseteq V$，八条规律中的$(i)(ii)(v)(vi)(vii)(viii)$自动成立，只需封闭性即可保证$(iii)(iv)$（由$0\boldsymbol{\alpha}=\boldsymbol{0}$、$(-1)\boldsymbol{\alpha}=-\boldsymbol{\alpha}$，零元素与负元素都在$L$中）。

由$V$中向量组$\boldsymbol{\alpha}_1,\dots,\boldsymbol{\alpha}_m$所生成的子空间为
\[
L=\{\boldsymbol{\gamma}\mid \boldsymbol{\gamma}=\lambda_1\boldsymbol{\alpha}_1+\dots+\lambda_m\boldsymbol{\alpha}_m,\ \lambda_i\in F\}
=\operatorname{Span}\{\boldsymbol{\alpha}_1,\dots,\boldsymbol{\alpha}_m\},
\]
它是包含$\boldsymbol{\alpha}_1,\dots,\boldsymbol{\alpha}_m$的最小子空间。

\keyformula{L\le V\ \text{（子空间）}\iff L\ne\varnothing\ \text{且}\ L\ \text{对加法、数乘封闭}}

\keyformula{\operatorname{Span}\{\boldsymbol{\alpha}_1,\dots,\boldsymbol{\alpha}_m\}=\Bigl\{\textstyle\sum_{i=1}^{m}\lambda_i\boldsymbol{\alpha}_i\Bigm|\lambda_i\in F\Bigr\},\qquad \dim\operatorname{Span}\{\boldsymbol{\alpha}_1,\dots,\boldsymbol{\alpha}_m\}=R(\boldsymbol{\alpha}_1,\dots,\boldsymbol{\alpha}_m)}

\section{维数、基与坐标}

\subsection{基与维数}
第4章中由线性运算建立的概念（线性组合、线性相关、最大无关组等）只涉及线性运算，因此对一般线性空间中的向量仍然适用，可直接引用。

\begin{definition}[基、维数]
在线性空间$V$中，如果存在$n$个向量$\boldsymbol{\alpha}_1,\boldsymbol{\alpha}_2,\dots,\boldsymbol{\alpha}_n$，满足
\begin{enumerate}
    \item[(i)] $\boldsymbol{\alpha}_1,\dots,\boldsymbol{\alpha}_n$线性无关；
    \item[(ii)] $V$中任一向量总可由$\boldsymbol{\alpha}_1,\dots,\boldsymbol{\alpha}_n$线性表示，
\end{enumerate}
那么$\boldsymbol{\alpha}_1,\dots,\boldsymbol{\alpha}_n$就称为线性空间$V$的一个\textbf{基}（基底），$n$称为线性空间$V$的\textbf{维数}，记作$\dim V=n$。维数为$n$的线性空间称为\textbf{$n$维线性空间}，记作$V_n$。
\end{definition}

只含一个零向量的线性空间没有基，规定其维数为$0$。$n$维与零维线性空间统称\textbf{有限维线性空间}；本书不讨论无限维线性空间。

由第4章的理论：若把线性空间看作向量组，则$V$的基就是该向量组的最大无关组，$V$的维数就是该向量组的秩。对$n$维线性空间$V_n$，若$\boldsymbol{\alpha}_1,\dots,\boldsymbol{\alpha}_n$为其一个基，则
\[
V_n=\operatorname{Span}\{\boldsymbol{\alpha}_1,\dots,\boldsymbol{\alpha}_n\},
\]
这就比较清楚地显示出$V_n$的构造。

\subsection{坐标}
\begin{definition}[坐标]
设$\boldsymbol{\alpha}_1,\dots,\boldsymbol{\alpha}_n$是线性空间$V_n$的一个基，对于任一向量$\boldsymbol{\alpha}\in V$，总有且仅有一组有序数$x_1,x_2,\dots,x_n\in F$，使
\[
\boldsymbol{\alpha}=x_1\boldsymbol{\alpha}_1+x_2\boldsymbol{\alpha}_2+\dots+x_n\boldsymbol{\alpha}_n,
\]
这组有序数就称为向量$\boldsymbol{\alpha}$在该基中的\textbf{坐标}，并记作$(x_1,x_2,\dots,x_n)$。
\end{definition}

取定一个基以后，$V_n$中的向量$\boldsymbol{\alpha}$与$\mathbb{R}^{n}$中的坐标向量$\boldsymbol{x}=(x_1,\dots,x_n)^{T}$之间就有一一对应的关系，且这个对应关系\textbf{保持线性运算}：
\[
\boldsymbol{\alpha}\leftrightarrow\boldsymbol{x},\quad \boldsymbol{\beta}\leftrightarrow\boldsymbol{y}
\ \Longrightarrow\ 
k\boldsymbol{\alpha}+l\boldsymbol{\beta}\leftrightarrow k\boldsymbol{x}+l\boldsymbol{y}.
\]
因此$V_n$与$\mathbb{R}^{n}$有相同的结构，称$V_n$与$\mathbb{R}^{n}$\textbf{同构}。一般地，两个线性空间之间若存在保持线性运算的一一对应，就说它们同构；维数相等的有限维线性空间都同构，即\textbf{线性空间的结构完全由它的维数决定}。

但要注意：$\mathbb{R}^{n}$中的内积概念在一般线性空间$V_n$中不一定有意义——同构只保证线性运算，不保证内积。

\subsection{求坐标的步骤}
\begin{enumerate}
    \item 设所求坐标为$\boldsymbol{x}=(x_1,\dots,x_n)^{T}$，按定义把$\boldsymbol{\alpha}=x_1\boldsymbol{\alpha}_1+\dots+x_n\boldsymbol{\alpha}_n$写成矩阵形式
    \[
    (\boldsymbol{\alpha}_1,\boldsymbol{\alpha}_2,\dots,\boldsymbol{\alpha}_n)\boldsymbol{x}=\boldsymbol{\alpha}.
    \]
    \item 若在$\mathbb{R}^{n}$中，令$A=(\boldsymbol{\alpha}_1,\dots,\boldsymbol{\alpha}_n)$，解方程组$A\boldsymbol{x}=\boldsymbol{\alpha}$（用初等行变换化$(A,\boldsymbol{\alpha})$为行最简形），所得解即坐标。
    \item 若题目只问"是否为基"，等价于判断$\boldsymbol{\alpha}_1,\dots,\boldsymbol{\alpha}_n$线性无关（$n$个向量在$n$维空间中），即$R(A)=n$或$|A|\ne0$。
\end{enumerate}

\corepoint{
（1）坐标是\textbf{相对于某个基}而言的，换基则坐标改变，写坐标必须说明"在基$\boldsymbol{\alpha}_1,\dots,\boldsymbol{\alpha}_n$下"。\\
（2）求坐标是解线性方程组，不是"把基按行排好直接读"；只有取自然基$\boldsymbol{e}_1,\dots,\boldsymbol{e}_n$时，坐标才恰好等于分量。\\
（3）多项式空间中求坐标，要先把向量用基"归并同类项"比较系数，本质仍是解方程组。
}

\section{基变换与坐标变换}

同一向量在不同基中的坐标不同，两者之间的关系由过渡矩阵刻画。

\subsection{基变换公式与过渡矩阵}
设$\boldsymbol{\alpha}_1,\dots,\boldsymbol{\alpha}_n$及$\boldsymbol{\beta}_1,\dots,\boldsymbol{\beta}_n$是线性空间$V_n$的两个基，把新基向量用旧基线性表示：
\[
\begin{cases}
\boldsymbol{\beta}_1=p_{11}\boldsymbol{\alpha}_1+p_{21}\boldsymbol{\alpha}_2+\dots+p_{n1}\boldsymbol{\alpha}_n,\\
\boldsymbol{\beta}_2=p_{12}\boldsymbol{\alpha}_1+p_{22}\boldsymbol{\alpha}_2+\dots+p_{n2}\boldsymbol{\alpha}_n,\\
\qquad\qquad\cdots\cdots\cdots\cdots\\
\boldsymbol{\beta}_n=p_{1n}\boldsymbol{\alpha}_1+p_{2n}\boldsymbol{\alpha}_2+\dots+p_{nn}\boldsymbol{\alpha}_n,
\end{cases}
\]
记$P=(p_{ij})_{n\times n}$，用矩阵形式写作
\[
(\boldsymbol{\beta}_1,\boldsymbol{\beta}_2,\dots,\boldsymbol{\beta}_n)
=(\boldsymbol{\alpha}_1,\boldsymbol{\alpha}_2,\dots,\boldsymbol{\alpha}_n)
\begin{pmatrix}p_{11}&p_{12}&\dots&p_{1n}\\p_{21}&p_{22}&\dots&p_{2n}\\\vdots&\vdots&\ddots&\vdots\\p_{n1}&p_{n2}&\dots&p_{nn}\end{pmatrix}
=(\boldsymbol{\alpha}_1,\dots,\boldsymbol{\alpha}_n)P .
\]
此式称为\textbf{基变换公式}，矩阵$P$称为由基$\boldsymbol{\alpha}_1,\dots,\boldsymbol{\alpha}_n$到基$\boldsymbol{\beta}_1,\dots,\boldsymbol{\beta}_n$的\textbf{过渡矩阵}。由于$\boldsymbol{\beta}_1,\dots,\boldsymbol{\beta}_n$线性无关，过渡矩阵$P$必可逆。

\subsection{坐标变换公式}
\begin{theorem}[坐标变换公式]
设$V_n$中的向量$\boldsymbol{\alpha}$在基$\boldsymbol{\alpha}_1,\dots,\boldsymbol{\alpha}_n$中的坐标为$(x_1,\dots,x_n)^{T}$，在基$\boldsymbol{\beta}_1,\dots,\boldsymbol{\beta}_n$中的坐标为$(x_1',\dots,x_n')^{T}$。若两个基满足基变换公式$(\boldsymbol{\beta}_1,\dots,\boldsymbol{\beta}_n)=(\boldsymbol{\alpha}_1,\dots,\boldsymbol{\alpha}_n)P$，则有
\[
\begin{pmatrix}x_1\\x_2\\\vdots\\x_n\end{pmatrix}
=P\begin{pmatrix}x_1'\\x_2'\\\vdots\\x_n'\end{pmatrix},
\qquad\text{或}\qquad
\begin{pmatrix}x_1'\\x_2'\\\vdots\\x_n'\end{pmatrix}
=P^{-1}\begin{pmatrix}x_1\\x_2\\\vdots\\x_n\end{pmatrix}.
\]
\end{theorem}

\textbf{证明}\quad 由$\boldsymbol{\alpha}=(\boldsymbol{\alpha}_1,\dots,\boldsymbol{\alpha}_n)\boldsymbol{x}=(\boldsymbol{\beta}_1,\dots,\boldsymbol{\beta}_n)\boldsymbol{x}'=(\boldsymbol{\alpha}_1,\dots,\boldsymbol{\alpha}_n)P\boldsymbol{x}'$，而$\boldsymbol{\alpha}_1,\dots,\boldsymbol{\alpha}_n$线性无关，故$\boldsymbol{x}=P\boldsymbol{x}'$。\hfill$\square$

定理的逆命题也成立：若任一向量的两种坐标都满足$\boldsymbol{x}=P\boldsymbol{x}'$，则两个基满足$(\boldsymbol{\beta})=(\boldsymbol{\alpha})P$。

\subsection{求过渡矩阵与坐标的方法}
\begin{enumerate}
    \item \textbf{代入法（定义法）}：把每个新基向量用旧基线性表示，解得系数，按"第$j$列是新基第$j$个向量的坐标"排成$P$。
    \item \textbf{初等变换法}：在$\mathbb{R}^{n}$中，令$A=(\boldsymbol{\alpha}_1,\dots,\boldsymbol{\alpha}_n)$，$B=(\boldsymbol{\beta}_1,\dots,\boldsymbol{\beta}_n)$，则$B=AP$，$P=A^{-1}B$。对$(A\mid B)$作初等行变换，当$A$化成$E$时，$B$就化成了$P$，即
    \[
    (A\mid B)\ \xrightarrow{\text{行变换}}\ (E\mid P).
    \]
    \item \textbf{坐标变换}：先用上法求$P$，再用$\boldsymbol{x}'=P^{-1}\boldsymbol{x}$（或解$P\boldsymbol{x}'=\boldsymbol{x}$）求新坐标。
\end{enumerate}

\keyformula{(\boldsymbol{\beta}_1,\dots,\boldsymbol{\beta}_n)=(\boldsymbol{\alpha}_1,\dots,\boldsymbol{\alpha}_n)P\quad(\text{基变换}),\qquad \boldsymbol{x}=P\boldsymbol{x}'\quad(\text{旧坐标}=\text{过渡矩阵}\times\text{新坐标})}

\keyformula{P=A^{-1}B,\qquad (A\mid B)\longrightarrow(E\mid P)}

\corepoint{
（1）方向极易记反：$P$是"由旧基\textbf{到}新基"的过渡矩阵，满足$(\text{新基})=(\text{旧基})P$；坐标变换则是$\boldsymbol{x}=P\boldsymbol{x}'$，即"旧坐标$=$过渡矩阵$\times$新坐标"，两者方向一致，不必再记$P^{-1}$在左边还是右边，直接用$\boldsymbol{x}=P\boldsymbol{x}'$推。\\
（2）$P$的列必须是新基向量在旧基下的坐标；若把$P$的行写成坐标，转置一下就错了。\\
（3）用$(A\mid B)\to(E\mid P)$时只能做\textbf{行变换}；若$A$不可逆说明$B$不是基，应回头检查。
}

\section{线性变换}

\subsection{映射与线性变换}
\begin{definition}[映射]
设有两个非空集合$A,B$，如果对于$A$中任一元素$\boldsymbol{\alpha}$，按照一定的规则，总有$B$中一个确定的元素$\boldsymbol{\beta}$和它对应，那么这个对应规则称为从$A$到$B$的\textbf{映射}，记作$T$，并记$\boldsymbol{\beta}=T(\boldsymbol{\alpha})$或$\boldsymbol{\beta}=T\boldsymbol{\alpha}$。$\boldsymbol{\beta}$称为$\boldsymbol{\alpha}$在映射$T$下的\textbf{像}，$\boldsymbol{\alpha}$称为$\boldsymbol{\beta}$的\textbf{原像}；$A$称为$T$的\textbf{定义域}，像的全体$T(A)=\{\boldsymbol{\beta}\mid\boldsymbol{\beta}=T(\boldsymbol{\alpha}),\ \boldsymbol{\alpha}\in A\}$称为\textbf{像集}。
\end{definition}

映射是函数概念的推广。

\begin{definition}[线性映射、线性变换]
设$V_n,U_m$分别是$n$维和$m$维线性空间，$T$是从$V_n$到$U_m$的映射。如果$T$满足
\begin{enumerate}
    \item[(i)] 任给$\boldsymbol{\alpha}_1,\boldsymbol{\alpha}_2\in V_n$，有$T(\boldsymbol{\alpha}_1+\boldsymbol{\alpha}_2)=T(\boldsymbol{\alpha}_1)+T(\boldsymbol{\alpha}_2)$；
    \item[(ii)] 任给$\boldsymbol{\alpha}\in V_n$，$\lambda\in F$，有$T(\lambda\boldsymbol{\alpha})=\lambda T(\boldsymbol{\alpha})$，
\end{enumerate}
那么$T$就称为从$V_n$到$U_m$的\textbf{线性映射}。$V_n$到自身的线性映射称为$V_n$中的\textbf{线性变换}。
\end{definition}

简言之，线性映射就是\textbf{保持线性组合}的映射：
\[
T(k_1\boldsymbol{\alpha}_1+\dots+k_m\boldsymbol{\alpha}_m)=k_1T(\boldsymbol{\alpha}_1)+\dots+k_mT(\boldsymbol{\alpha}_m).
\]

典型例子：多项式空间的求导运算$D$（$Dp=p'$）是线性变换；$\mathbb{R}^{n}$中由$n$阶矩阵$A$确定的变换$T(\boldsymbol{x})=A\boldsymbol{x}$是线性变换；平面上的旋转变换、把向量投影到某平面上的变换都是线性变换。而$T(p)=1$（把任何多项式映成常数$1$）不是线性变换，因为$T(p+q)=1\ne1+1=T(p)+T(q)$。

零变换$\Gamma_0$（把一切向量映成$\boldsymbol{0}$）与恒等变换$I$（$I(\boldsymbol{\alpha})=\boldsymbol{\alpha}$）都是线性变换。

\subsection{线性变换的基本性质}
设$T$是$V_n$中的线性变换，则
\begin{enumerate}
    \item $T(\boldsymbol{0})=\boldsymbol{0}$，$T(-\boldsymbol{\alpha})=-T(\boldsymbol{\alpha})$；
    \item 若$\boldsymbol{\beta}=k_1\boldsymbol{\alpha}_1+\dots+k_m\boldsymbol{\alpha}_m$，则$T(\boldsymbol{\beta})=k_1T(\boldsymbol{\alpha}_1)+\dots+k_mT(\boldsymbol{\alpha}_m)$（线性变换由它在一组基上的像唯一确定）；
    \item 若$\boldsymbol{\alpha}_1,\dots,\boldsymbol{\alpha}_m$线性相关，则$T(\boldsymbol{\alpha}_1),\dots,T(\boldsymbol{\alpha}_m)$也线性相关；
    \item 像集$T(V_n)$是$V_n$的子空间，称为$T$的\textbf{像空间}；
    \item 使$T(\boldsymbol{\alpha})=\boldsymbol{0}$的$\boldsymbol{\alpha}$的全体
    \[
    N_T=\ker T=\{\boldsymbol{\alpha}\in V_n\mid T(\boldsymbol{\alpha})=\boldsymbol{0}\}
    \]
    也是$V_n$的子空间，称为$T$的\textbf{核}。
\end{enumerate}

注意性质3的\textbf{逆命题不成立}：线性无关的向量组经线性变换后可能变得线性相关（例如投影变换把平面内两条不共线向量都投到同一条直线上）。这是因为线性变换可能"压平"空间。

\corepoint{
（1）判断一个映射是否为线性变换，只需验加法与数乘两条；常用来否定的是"常数项不为零"的映射（如$T(\boldsymbol{x})=A\boldsymbol{x}+\boldsymbol{b}$，$\boldsymbol{b}\ne\boldsymbol{0}$不是线性变换）。\\
（2）$T(\boldsymbol{0})=\boldsymbol{0}$是必要条件：若$T(\boldsymbol{0})\ne\boldsymbol{0}$，则$T$一定不是线性变换。\\
（3）像空间与核都是子空间，不要把"核"与"零向量"混淆：核一般含无穷多向量，只是其中包含零向量。
}

\section{线性变换的矩阵表示式}

\subsection{线性变换在给定基下的矩阵}
设$T$是$V_n$中的线性变换，$\boldsymbol{\varepsilon}_1,\boldsymbol{\varepsilon}_2,\dots,\boldsymbol{\varepsilon}_n$是$V_n$的一个基。由于$T(\boldsymbol{\varepsilon}_j)$仍在$V_n$中，故可由该基唯一地线性表示：
\[
\begin{cases}
T(\boldsymbol{\varepsilon}_1)=a_{11}\boldsymbol{\varepsilon}_1+a_{21}\boldsymbol{\varepsilon}_2+\dots+a_{n1}\boldsymbol{\varepsilon}_n,\\
T(\boldsymbol{\varepsilon}_2)=a_{12}\boldsymbol{\varepsilon}_1+a_{22}\boldsymbol{\varepsilon}_2+\dots+a_{n2}\boldsymbol{\varepsilon}_n,\\
\qquad\qquad\cdots\cdots\cdots\cdots\\
T(\boldsymbol{\varepsilon}_n)=a_{1n}\boldsymbol{\varepsilon}_1+a_{2n}\boldsymbol{\varepsilon}_2+\dots+a_{nn}\boldsymbol{\varepsilon}_n .
\end{cases}
\]
记
\[
A=\begin{pmatrix}
a_{11}&a_{12}&\dots&a_{1n}\\
a_{21}&a_{22}&\dots&a_{2n}\\
\vdots&\vdots&\ddots&\vdots\\
a_{n1}&a_{n2}&\dots&a_{nn}
\end{pmatrix},
\]
则上式可形式地写成
\[
T(\boldsymbol{\varepsilon}_1,\boldsymbol{\varepsilon}_2,\dots,\boldsymbol{\varepsilon}_n)
=\bigl(T(\boldsymbol{\varepsilon}_1),T(\boldsymbol{\varepsilon}_2),\dots,T(\boldsymbol{\varepsilon}_n)\bigr)
=(\boldsymbol{\varepsilon}_1,\boldsymbol{\varepsilon}_2,\dots,\boldsymbol{\varepsilon}_n)A .
\]

\begin{definition}[线性变换的矩阵]
上述矩阵$A$称为线性变换$T$在基$\boldsymbol{\varepsilon}_1,\boldsymbol{\varepsilon}_2,\dots,\boldsymbol{\varepsilon}_n$下的\textbf{矩阵}。
\end{definition}

矩阵$A$由基的像$T(\boldsymbol{\varepsilon}_1),\dots,T(\boldsymbol{\varepsilon}_n)$唯一确定；反之，给定一个$n$阶矩阵$A$，也就给定了基的像，从而唯一确定一个以$A$为矩阵的线性变换。因此在\textbf{取定一组基之后，线性变换与$n$阶矩阵一一对应}。

若$V_n$中向量$\boldsymbol{\alpha}$在该基下的坐标为$\boldsymbol{x}=(x_1,\dots,x_n)^{T}$，则$T(\boldsymbol{\alpha})$在同一基下的坐标为
\[
\boldsymbol{y}=A\boldsymbol{x},
\quad\text{即}\quad
T\bigl((\boldsymbol{\varepsilon}_1,\dots,\boldsymbol{\varepsilon}_n)\boldsymbol{x}\bigr)
=(\boldsymbol{\varepsilon}_1,\dots,\boldsymbol{\varepsilon}_n)A\boldsymbol{x}.
\]

\subsection{不同基下矩阵的关系}
\begin{theorem}[同一线性变换在不同基下的矩阵相似]
设线性空间$V_n$中取定两个基$\boldsymbol{\alpha}_1,\dots,\boldsymbol{\alpha}_n$和$\boldsymbol{\beta}_1,\dots,\boldsymbol{\beta}_n$，由基$\boldsymbol{\alpha}_1,\dots,\boldsymbol{\alpha}_n$到基$\boldsymbol{\beta}_1,\dots,\boldsymbol{\beta}_n$的过渡矩阵为$P$；$V_n$中的线性变换$T$在这两个基下的矩阵依次为$A$和$B$，则
\[
B=P^{-1}AP .
\]
\end{theorem}

\textbf{证明}\quad 由$(\boldsymbol{\beta}_1,\dots,\boldsymbol{\beta}_n)=(\boldsymbol{\alpha}_1,\dots,\boldsymbol{\alpha}_n)P$得
\[
\begin{aligned}
T(\boldsymbol{\beta}_1,\dots,\boldsymbol{\beta}_n)
&=T\bigl((\boldsymbol{\alpha}_1,\dots,\boldsymbol{\alpha}_n)P\bigr)
=\bigl(T(\boldsymbol{\alpha}_1,\dots,\boldsymbol{\alpha}_n)\bigr)P\\
&=(\boldsymbol{\alpha}_1,\dots,\boldsymbol{\alpha}_n)AP
=(\boldsymbol{\beta}_1,\dots,\boldsymbol{\beta}_n)P^{-1}AP,
\end{aligned}
\]
又$T(\boldsymbol{\beta}_1,\dots,\boldsymbol{\beta}_n)=(\boldsymbol{\beta}_1,\dots,\boldsymbol{\beta}_n)B$，而$\boldsymbol{\beta}_1,\dots,\boldsymbol{\beta}_n$线性无关，故$B=P^{-1}AP$。\hfill$\square$

定理表明：\textbf{同一线性变换在不同基下的矩阵是相似的，相似变换矩阵就是过渡矩阵$P$}。这是"相似"概念的几何意义——相似矩阵正是同一个线性变换在不同基下的矩阵。

\subsection{线性变换的秩}
\begin{definition}[线性变换的秩]
线性变换$T$的像空间$T(V_n)$的维数，称为线性变换$T$的\textbf{秩}。
\end{definition}

若$A$是$T$的矩阵，则$T$的秩就等于$R(A)$；若$R(A)=r$，则$T$的核的维数为
\[
\dim\ker T=n-r .
\]
（因为$\ker T$就是齐次方程组$A\boldsymbol{x}=\boldsymbol{0}$的解空间。）

\keyformula{T(\boldsymbol{\varepsilon}_1,\dots,\boldsymbol{\varepsilon}_n)=(\boldsymbol{\varepsilon}_1,\dots,\boldsymbol{\varepsilon}_n)A,\qquad T(\boldsymbol{\alpha})\ \text{的坐标}=A\boldsymbol{x}}

\keyformula{B=P^{-1}AP\ (\text{换基}),\qquad \operatorname{rank}T=R(A),\qquad \dim\ker T=n-R(A)}

\corepoint{
（1）求$T$在基$\boldsymbol{\varepsilon}_1,\dots,\boldsymbol{\varepsilon}_n$下的矩阵，关键是算$T(\boldsymbol{\varepsilon}_j)$并把它们\textbf{用该基}表示，系数按\textbf{列}排成$A$；若把系数按行排，得到的是$A^{T}$。\\
（2）同一变换换基后矩阵变为$P^{-1}AP$，不是$P^{T}AP$（那是合同变换），也不是$P^{-1}AP^{-1}$。\\
（3）"线性变换与矩阵一一对应"依赖于\textbf{取定一组基}；换基后对应同一个变换的矩阵不同，但它们相似。\\
（4）$\dim$像$+\dim$核$=n$是验算的有力工具。
}
